\section{Conclusiones y recomendaciones}
\label{sec:conclusiones}

\subsection{Conclusiones}

\begin{enumerate}
    \item \textbf{Detección y seguimiento local.} YOLO26m con el seguidor
    híbrido por cámara siguió a las 16 personas de la primera prueba; el
    ajuste de umbrales bajó los identificadores locales de 88 a 80 y los
    de personas lejanas de 7 a 1, de modo que una persona aún puede
    partirse en varios tramos dentro de una cámara. El modelo procesa unos
    17 cuadros por segundo por video con dos cámaras 1080p en la RTX~4060
    Laptop (Secciones~\ref{sec:pruebas} y~\ref{sec:validacion-esan}).
    \item \textbf{Re-identificación sin rostro.} La apariencia corporal
    bastó para asociar personas entre cámaras con una precisión de 0.99 y
    un \textit{recall} de 0.86 en la primera prueba; en el comedor cada
    persona recibió el mismo identificador en los dos teléfonos, y en las
    grabaciones del 9 de noviembre se siguió a 149 personas en el patio
    (41 en las dos cámaras) y a 128 en el pasillo (20 en las dos)
    (Sección~\ref{sec:mas-gente}).
    \item \textbf{Plano 2D.} La homografía ubicó en el plano el 93\,\% de
    las detecciones con un error de calibración mediano de 0.28 a 0.48~m,
    y las referencias del en vivo ubicaron grabaciones sin modelo de
    cámara (24 de 24 cuadros reconocidos).
    \item \textbf{Patrones históricos.} KDE, PrefixSpan y el grafo
    origen-destino generan mapas de calor, rutas frecuentes y flujos entre
    zonas a partir de las sesiones guardadas (Sección~\ref{sec:algoritmos}).
    \item \textbf{Métricas y tablero.} La plataforma muestra en vivo y por
    sesión los conteos, permanencias, mapas de calor y rutas, exporta el
    reporte en PDF y controla el acceso por roles, con permisos de
    exportación aprobados por el superadmin (Sección~\ref{sec:arquitectura}).
    \item \textbf{Datos.} PostgreSQL/PostGIS conserva la trazabilidad de
    cámara, identidad seudónima, trayectoria y evento sin imágenes. La
    memoria de identidades del en vivo sí guarda vectores de apariencia,
    que se borran a las 24~horas sin ver a la persona; los puntos crudos se
    borran a los 90 días y quedan solo los agregados.
    \item \textbf{Salvaguardas.} El sistema no usa el rostro ni estima el
    género, y la retención es corta; quedan abiertas las brechas legales
    de la Sección~\ref{sec:analisis-legal}, en especial el tratamiento de
    menores y la calificación de los vectores de Re-ID como dato
    biométrico.
    \item \textbf{Costos y negocio.} Procesar en una computadora propia
    cuesta unos US\$~70 al mes, frente a US\$~379--389 de una T4 a demanda
    en las grandes nubes (Cuadro~\ref{tab:nube-proveedores}); con ese costo,
    el plan de negocio da un VAN de S/~798\,320 al 20\,\% y sigue positivo
    con los flujos reducidos a la cuarta parte (Cuadro~\ref{tab:van-sens}).
\end{enumerate}

\subsection{Recomendaciones}

\begin{itemize}
    \item Validar el sistema con el CCTV real del aeropuerto, en un piloto
    acordado con LAP, porque las pruebas se hicieron en ESAN con tres
    cámaras fijas y con teléfonos.
    \item Cerrar antes de cualquier despliegue las decisiones legales
    pendientes: opinión legal sobre los vectores de Re-ID, tratamiento de
    menores, contrato de encargo con el cliente y registro del banco de
    datos (Cuadro~\ref{tab:matriz-legal}).
    \item Medir el sesgo del detector y del Re-ID con grupos diversos de
    personas, como propone la matriz de cumplimiento~\cite{matrices2026}.
    \item Cotizar la computadora con GPU, que hoy es un supuesto de
    S/~4\,500, y elegir entre computadora propia y nube según las horas de
    uso al mes (Cuadro~\ref{tab:nube-proveedores}).
\end{itemize}
